% =========================================================================
% Seguimiento de riesgos — 18-10-2026 (S8, segundo sprint de F2).
% Revisión por cadencia de los críticos (R1–R7 y R10); R8 por cambio de
% medida (semana a cero); R11 por disparador observado el 15-10; R12 por
% cambio de medida (fin de las entregas quincenales) y vigilancia de su
% indicador; R14 por seguimiento del indicador que saltó en S7. R15 no se
% adelanta (el límite de sesión del 15-10 no es su indicador): toca por
% cadencia en S9. Auditado contra el dev-journal (12-10,
% 15-10, 17-10 y 18-10) y git antes de emitir la v7 del plan. Documento
% independiente de la memoria; creado copiando TEMPLATE.tex (ver README.md
% de esta carpeta).
% Compilar DESDE ESTA CARPETA (2 pases, sin biber):
%   pdflatex --enable-installer -interaction=nonstopmode -file-line-error
%     -output-directory=build 2026-10-18-S8.tex
% =========================================================================
\documentclass[11pt,a4paper]{article}

% =========================================================================
% Preámbulo COMPARTIDO de los documentos independientes de riesgos
% (docs/thesis/plan-riesgos/ y docs/thesis/seguimiento-riesgos/).
%
% Estos documentos NO forman parte de main.tex (la memoria): son entregables
% sueltos en clase article, con el diseño heredado de los .tex de Gestión de
% Proyectos del autor, pero adaptados a la configuración global de la
% memoria (Carlito, márgenes 2,5 cm, parskip 6 pt, títulos de la guía).
% Un solo preámbulo para toda la familia: un ajuste de estilo aquí se
% propaga a todos los seguimientos futuros (uno por sprint hasta S37).
%
% RUTAS: los \input son relativos a la carpeta del documento, no a esta;
% compilar SIEMPRE desde la carpeta del .tex (plan-riesgos/ o
% seguimiento-riesgos/):
%   pdflatex --enable-installer -interaction=nonstopmode -file-line-error
%     -output-directory=build <documento>.tex        (2 pases; sin biber)
% =========================================================================

% --- Codificación (idéntico a packages.tex de la memoria) ---
\usepackage[utf8]{inputenc}
\usepackage[T1]{fontenc}

% --- Idioma: solo español (aquí no hay resumen trilingüe) ---
% es-noshorthands por coherencia con la memoria: acentos en UTF-8, sin
% caracteres activos (el conflicto con biblatex-apa no aplica aquí, pero
% mantener una única configuración de babel evita sorpresas al copiar texto
% entre la memoria y estos documentos).
\usepackage[spanish, es-tabla, es-noshorthands]{babel}

% --- Márgenes: 2,5 cm laterales e inferior, como la memoria ---
% El margen superior se reparte para alojar la cabecera enmarcada del
% diseño de referencia (2 cm de caja + separación) sin comerse el cuerpo:
% el texto arranca a 4 cm del borde y el ancho de texto sigue siendo 16 cm.
\usepackage[a4paper, left=2.5cm, right=2.5cm, bottom=2.5cm,
            top=1.2cm, includehead, headheight=2.2cm, headsep=0.6cm]{geometry}

% --- Tipografía: Carlito, la misma decisión que en style.tex ---
% (clon métrico libre de Calibri; pdflatex no puede usar la real).
% No se replica el arreglo de versalitas de style.tex: aquí no hay
% bibliografía APA y ningún texto usa \textsc.
\usepackage[sfdefault]{carlito}

% --- Espaciado: interlineado sencillo + 6 pt entre párrafos (guía) ---
\usepackage[skip=6pt]{parskip}

% --- Tablas ---
% longtable: el registro y el plan de respuesta superan una página.
% AVISO ESPERADO: "ignored error: Infinite glue shrinkage" en cada salto de
% página de una longtable (bug conocido e inofensivo, latex3/latex2e#1907;
% ver el comentario largo en packages.tex de la memoria).
\usepackage{array}
\usepackage{longtable}
\usepackage{tabularx}
% \RaggedRight (ragged2e) en vez de \raggedright en columnas muy estrechas:
% el \raggedright de LaTeX desactiva de facto la partición silábica y una
% palabra larga («Identificado») desborda una columna de 1,2 cm; la
% variante de ragged2e mantiene el guionado.
\usepackage{ragged2e}
% Las columnas X se componen como m{...}: el contenido de cada celda queda
% centrado verticalmente (aspecto de formulario del diseño de referencia).
\renewcommand{\tabularxcolumn}[1]{m{#1}}
% Tipos de columna del diseño de referencia:
%   L/C = ancho fijo alineado izquierda/centro; Y/Z = X alineada izq./centro.
% En columnas estrechas se prefiere Y (raggedright) a X justificada: la
% justificación en 3-4 cm abre ríos de blancos.
\newcolumntype{L}[1]{>{\raggedright\arraybackslash}m{#1}}
\newcolumntype{C}[1]{>{\centering\arraybackslash}m{#1}}
\newcolumntype{Y}{>{\raggedright\arraybackslash}X}
\newcolumntype{Z}{>{\centering\arraybackslash}X}
% Filas con aire en todas las tablas (el diseño de referencia lo fijaba
% tabla a tabla; aquí una sola vez).
\renewcommand{\arraystretch}{1.25}

% --- Utilidades ---
\usepackage{xcolor}    % \markfield (resaltado) y \todonote (macros.tex)
\usepackage{amssymb}   % \square, \blacksquare, \circlearrowright
\usepackage{xspace}    % lo necesita \bukit (macros.tex)
\usepackage{enumitem}
\setlist[itemize]{leftmargin=*, topsep=2pt, itemsep=2pt}
\setlist[enumerate]{leftmargin=*, topsep=2pt, itemsep=2pt}

% --- Títulos según la guía de estilo, como en style.tex de la memoria ---
% (sin \chapter: la clase article no lo tiene; el resto, idéntico).
\usepackage{titlesec}
\titleformat{\section}
  {\normalfont\bfseries\fontsize{12pt}{14.4pt}\selectfont\raggedright}
  {\thesection.}{0.5em}{\MakeUppercase}
\titleformat{\subsection}
  {\normalfont\bfseries\fontsize{12pt}{14.4pt}\selectfont}
  {\thesubsection.}{0.5em}{}
\titleformat{\subsubsection}
  {\normalfont\itshape\fontsize{12pt}{14.4pt}\selectfont}
  {\thesubsubsection.}{0.5em}{}

% --- Paginación dinámica (regla estricta del autor, 26-07-2026) ---
% PROHIBIDO el salto de página manual también en estos documentos: el
% espacio lo garantizan estos ganchos (o un \Needspace* local ante un bloque
% indivisible), nunca un \newpage. Los .tex de referencia sí usaban
% \newpage; es la única desviación deliberada respecto a ellos.
\usepackage{needspace}
\AddToHook{cmd/section/before}{\Needspace*{5\baselineskip}}
\AddToHook{cmd/subsection/before}{\Needspace*{5\baselineskip}}
\AddToHook{cmd/subsubsection/before}{\Needspace*{4\baselineskip}}
\widowpenalty=10000
\clubpenalty=10000
\displaywidowpenalty=10000
\raggedbottom
% Estos formularios componen mucho texto en columnas de 2-4 cm: se concede
% a TeX un margen de estirado de emergencia antes de desbordar la línea
% (evita Overfull \hbox menores sin tocar el contenido).
\setlength{\emergencystretch}{1em}

% --- Datos globales y macros de la memoria (única fuente: no duplicar) ---
% =========================================================================
% Variables globales del documento.
% Cambiar un dato AQUÍ lo actualiza en todo el documento (portada,
% metadatos del PDF...). No escribir nunca estos datos "a mano" en otros
% archivos: usar siempre estos comandos.
% =========================================================================

% --- Trabajo ---
\newcommand{\thesisTitle}{BUKIT: red social orientada a objetivos}% TODO: título definitivo
\newcommand{\thesisSubtitle}{Memoria del Trabajo de Fin de Grado}%
\newcommand{\thesisDegree}{Grado en Ingeniería Informática de Gestión y Sistemas de Información}%
\newcommand{\thesisCourse}{2026--2027}% curso académico
\newcommand{\thesisDate}{Bilbao, \today}% TODO: <Lugar, día, mes, año> de la entrega/defensa
\newcommand{\thesisYear}{2027}% año de elaboración y defensa

% --- Personas (la plantilla del centro usa "Apellido1 Apellido2, Nombre") ---
\newcommand{\thesisAuthor}{Liviu Deleanu}% nombre natural (para metadatos del PDF)
\newcommand{\thesisAuthorFormal}{Deleanu, Liviu}% TODO: confirmar formato con dos apellidos
\newcommand{\thesisDirector}{Villamañe Gironés, Mikel}%
\newcommand{\thesisCoDirector}{Apellido1 Apellido2, Nombre}% TODO: solo si hay codirector/a

% --- Universidad ---
\newcommand{\thesisUniversity}{Universidad del País Vasco / Euskal Herriko Unibertsitatea}%
\newcommand{\thesisFaculty}{Escuela de Ingeniería de Bilbao}%

% --- Palabras clave del resumen trilingüe ---
\newcommand{\thesisKeywordsEs}{red social, objetivos, monorepo, TypeScript, PostgreSQL}% TODO
\newcommand{\thesisKeywordsEu}{sare soziala, helburuak, monorepo, TypeScript, PostgreSQL}% TODO: revisar euskera
\newcommand{\thesisKeywordsEn}{social network, goals, monorepo, TypeScript, PostgreSQL}% TODO

% =========================================================================
% Comandos propios. Definir aquí cualquier cosa que se repita en el texto:
% si mañana cambia cómo se escribe, se cambia en un único sitio.
% =========================================================================

% Nombre del producto, siempre con el mismo formato.
\newcommand{\bukit}{\textsc{Bukit}\xspace}

% Términos en inglés, en cursiva de forma consistente: \eng{backend}
\newcommand{\eng}[1]{\textit{#1}}

% Notas de trabajo, visibles en rojo en el PDF. Buscar "todonote" y
% eliminarlas todas antes de entregar.
\newcommand{\todonote}[1]{\textcolor{red}{\textbf{[TODO: #1]}}}


% --- Cabecera y pie (diseño de referencia, sin logos) ---
% La caja de cabecera replica la del .tex de referencia (celda ancha de
% contexto + celda estrecha de identidad). La celda derecha llevaba el logo
% del centro: aquí va texto porque los logos UPV/EHU siguen pendientes de
% la fase de diseño (ver style.tex); cuando existan, sustituir por
% \includegraphics. 11 + 4,9 = 15,9 cm < 16 cm de texto: sin desbordes.
\usepackage{fancyhdr}
\newcommand{\riesgosHeader}{%
  \setlength{\tabcolsep}{0pt}%
  \renewcommand{\arraystretch}{1}%
  \begin{tabular}{|p{11cm}|p{4.9cm}|}
    \hline
    \parbox[c][2cm][c]{11cm}{\centering\bfseries
      {\fontsize{16}{19}\selectfont TRABAJO DE FIN DE GRADO\\[4pt]Curso \thesisCourse}%
    } &
    \parbox[c][2cm][c]{4.9cm}{\centering
      {\fontsize{14}{17}\selectfont\bfseries\bukit}\\[3pt]
      {\small\thesisFaculty}%
    }\\
    \hline
  \end{tabular}%
}
\pagestyle{fancy}
\fancyhf{}
\renewcommand{\headrulewidth}{0pt}
\fancyhead[C]{\riesgosHeader}
\fancyfoot[C]{\thepage}
% La primera página usa estilo "plain": se redefine para que también lleve
% la cabecera (mismo truco que el .tex de referencia).
\fancypagestyle{plain}{%
  \fancyhf{}%
  \renewcommand{\headrulewidth}{0pt}%
  \fancyhead[C]{\riesgosHeader}%
  \fancyfoot[C]{\thepage}%
}

% --- Símbolos y campos de formulario del diseño de referencia ---
\newcommand{\fieldline}[1]{\underline{\hspace{#1}}}
\newcommand{\bigline}{\noindent\rule{\textwidth}{0.4pt}}
\newcommand{\checkbox}{\(\square\)}
\newcommand{\checkedbox}{\(\blacksquare\)}
% Símbolo ⟳ = "si se rellena este apartado, nueva versión del plan".
% El original lo resolvía con un SVG externo (paquete svg + Inkscape);
% aquí basta el símbolo de amssymb: cero dependencias nuevas.
\newcommand{\updatesymbol}{\raisebox{-0.2ex}{\(\circlearrowright\)}}
% Resaltado amarillo del identificador de riesgo en los títulos.
\newcommand{\markfield}[1]{\colorbox{yellow!35}{\strut #1}}

% --- Enlaces y metadatos del PDF (al final, como exige hyperref) ---
% El pdftitle lo fija cada documento con \hypersetup tras cargar esto.
\usepackage[hidelinks]{hyperref}
\hypersetup{pdfauthor={\thesisAuthor}}

\hypersetup{pdftitle={Seguimiento de riesgos — 18-10-2026 (S8)}}

\begin{document}

\begin{center}
  {\LARGE\bfseries Seguimiento de riesgos — S8}\\[6pt]
\end{center}

\textbf{Fecha:} 18-10-2026 · \textbf{Sprint:} S8 (12-10 \(\to\) 18-10-2026, segundo sprint de F2) · \textbf{Responsable:} \thesisAuthor

Regla de revisión: \textbf{críticos} cada sprint; \textbf{significativos} cada 2 sprints; \textbf{moderados} cada 4 sprints; cualquier riesgo, \textbf{de inmediato}, si se observa su indicador clave.
Una semana declarada a cero (cap.~10 de la memoria) no cierra sprint y no tiene seguimiento: lo que tocara por cadencia se revisa en el siguiente.

\textbf{El símbolo \updatesymbol\ indica que, si se rellena ese apartado, hay que emitir una nueva versión completa del plan de riesgos} en \texttt{../plan-riesgos/} y registrar el cambio en la tabla de seguimiento del cap.~7 de la memoria.

Riesgos revisados:
\textbf{R1--R7 y R10 por cadencia} (críticos, semanal);
\textbf{R8 por cambio de medida} (la semana a cero, decisión del autor del 17-10);
\textbf{R11 por disparador} observado el 15-10;
\textbf{R12 por cambio de medida} (fin de las entregas quincenales) y vigilancia de su indicador;
y \textbf{R14 por seguimiento} del indicador que saltó en S7 (activo en la v6).
R15 no se adelanta: el límite de sesión agotado el 15-10 no es su indicador, que mide la cuota semanal y las tareas recurrentes; queda anotado en «Otros eventos relevantes» y toca por cadencia en S9.
Próximas revisiones por cadencia: significativos (R8, R9, R11, R14 y R15) en S9 y moderados (R12 y R13) en S11.

\textbf{Plan emitido a partir de este seguimiento:} versión 7, con apartados \updatesymbol\ rellenados en R1, R3, R8, R10 y R12 (medidas e indicadores) y el estado de R11 y R14.\newline
Archivo: \path{../plan-riesgos/plan-riesgos-v7-18-10-2026.tex}.

Datos del sprint: velocidad de S8 = \textbf{2 SP} (US-07, cerrada el 17-10 y verificada en producción; con ella se completa EP02); horas reales según el diario, 8~h~30 (30~min el 12-10, 2~h~30 el 15-10, 4~h el 17-10 y 1~h~30 el 18-10), frente a las 5~h declaradas: 170~\%.
Sesiones el lunes, el jueves, el sábado y el domingo.
Equivalencia observada: no separable, porque las horas del sprint se repartieron entre la historia, el refactor de la API, los cambios del método y los documentos de riesgos, y el diario no las desglosa por tarea; se mantiene la de planificación, 2,0~h/SP, y la decisión sigue fijada para S9 (acuerdo de S7).

\bigline

\section{Medidas de mitigación (prevención)}

Para cada medida de los riesgos que tocan revisión en este seguimiento: ¿se está aplicando?
Si no, justificar y decidir si se mantiene/modifica/elimina.

% Bloques de 3-4 filas (tabularx no parte entre páginas; receta del S2).
\Needspace*{24\baselineskip}
\underline{\textbf{Riesgos críticos (R1--R4)}}

\begin{tabularx}{\textwidth}{|C{0.9cm}|Y|Y|C{2.2cm}|L{2.4cm}|}
\hline
\multicolumn{1}{|C{0.9cm}|}{\textbf{Id}} &
\multicolumn{1}{>{\centering\arraybackslash}X|}{\textbf{Medida de mitigación}} &
\multicolumn{1}{>{\centering\arraybackslash}X|}{\textbf{¿Aplicada?}} &
\multicolumn{1}{C{2.2cm}|}{\textbf{Justificación (del NO)}} &
\multicolumn{1}{C{2.4cm}|}{\textbf{Decisión}} \\
\hline
R1 & Toda idea nueva entra por el backlog. & \checkedbox\ Sí (el sprint comprometió solo la US-07; el refactor de la API no añade funciones, y EP04 y EP05 se refinaron en el backlog, no en el sprint).\newline \checkbox\ No & — & \checkedbox\ Mantener (se reescribe el indicador, abajo)\newline \checkbox\ Modificar\newline \checkbox\ Eliminar \\
\hline
R2 & Medir velocidad y recalibrar en cada cierre. & \checkedbox\ Sí (velocidad medida, 2 SP; la recalibración sigue fijada para S9).\newline \checkbox\ No & — & \checkedbox\ Mantener\newline \checkbox\ Modificar\newline \checkbox\ Eliminar \\
\hline
R3 & Diario al cierre de cada sesión; memoria con hueco; ninguna historia cerrada sin documentar; auditoría semanal de coherencia. & \checkedbox\ Sí (cap.~9 actualizado el mismo 17-10 con la US-07 y el refactor; EP04 y EP05 refinadas en el cap.~10 el 15-10; informe de coherencia del 11-10 aplicado el 15-10: 19 de sus 23 hallazgos en la primera sesión, 3 en la segunda y 1 que ya no aplicaba).\newline \checkbox\ No & — & \checkbox\ Mantener\newline \checkedbox\ Modificar (cómo se aplica el informe, abajo)\newline \checkbox\ Eliminar \\
\hline
R4 & Citar al escribir; pendientes con nota visible. & \checkedbox\ Sí (las cinco fuentes nuevas de la US-07 se citaron al escribir los caps.~6 y 9).\newline \checkbox\ No & — & \checkedbox\ Mantener\newline \checkbox\ Modificar\newline \checkbox\ Eliminar \\
\hline
\end{tabularx}

\Needspace*{24\baselineskip}
\underline{\textbf{Riesgos críticos (R5--R7 y R10)}}

\begin{tabularx}{\textwidth}{|C{0.9cm}|Y|Y|C{2.2cm}|L{2.4cm}|}
\hline
\multicolumn{1}{|C{0.9cm}|}{\textbf{Id}} &
\multicolumn{1}{>{\centering\arraybackslash}X|}{\textbf{Medida de mitigación}} &
\multicolumn{1}{>{\centering\arraybackslash}X|}{\textbf{¿Aplicada?}} &
\multicolumn{1}{C{2.2cm}|}{\textbf{Justificación (del NO)}} &
\multicolumn{1}{C{2.4cm}|}{\textbf{Decisión}} \\
\hline
R5 & Revisar todo antes de integrar; nada sin comprobaciones. & \checkedbox\ Sí (el refactor se aplicó por partes, con doble comprobación, y el autor exigió la justificación de cada decisión antes de aprobarla; lint, tipos y 81 pruebas en verde; segunda pasada sobre los comentarios del código).\newline \checkbox\ No & — & \checkedbox\ Mantener\newline \checkbox\ Modificar\newline \checkbox\ Eliminar \\
\hline
R6 & Investigación acotada ante tecnología nueva. & \checkedbox\ Sí (la biblioteca de limitación de peticiones se eligió el 15-10 entre dos candidatas, con su documentación, y la historia se cerró en su sprint).\newline \checkbox\ No & — & \checkedbox\ Mantener\newline \checkbox\ Modificar\newline \checkbox\ Eliminar \\
\hline
R7 & Minimización desde el diseño (RGPD). & \checkedbox\ Sí (la US-07 maneja la IP del cliente: solo en contadores en memoria de quince minutos, agrupada por /56 en IPv6 y censurada en los logs junto con las cabeceras de proxy).\newline \checkbox\ No & — & \checkedbox\ Mantener\newline \checkbox\ Modificar\newline \checkbox\ Eliminar \\
\hline
R10 & Margen deliberado + documentación continua. & \checkedbox\ Sí (sin parón; cuatro sesiones repartidas en la semana).\newline \checkbox\ No & — & \checkedbox\ Mantener (se redefinen su indicador y su contingencia, abajo)\newline \checkbox\ Modificar\newline \checkbox\ Eliminar \\
\hline
\end{tabularx}

% Dos bloques: con las justificaciones de R11 y R12, las cuatro filas no caben
% en una página con su título.
\Needspace*{22\baselineskip}
\underline{\textbf{Riesgos significativos (R8 y R14)}}

\begin{tabularx}{\textwidth}{|C{0.9cm}|Y|Y|C{2.2cm}|L{2.4cm}|}
\hline
\multicolumn{1}{|C{0.9cm}|}{\textbf{Id}} &
\multicolumn{1}{>{\centering\arraybackslash}X|}{\textbf{Medida de mitigación}} &
\multicolumn{1}{>{\centering\arraybackslash}X|}{\textbf{¿Aplicada?}} &
\multicolumn{1}{C{2.2cm}|}{\textbf{Justificación (del NO)}} &
\multicolumn{1}{C{2.4cm}|}{\textbf{Decisión}} \\
\hline
R8 & Horas por sesión en el diario; cada retro compara. & \checkedbox\ Sí (8~h~30 registradas frente a 5 declaradas).\newline \checkbox\ No & — & \checkedbox\ Mantener (se añade la semana a cero, abajo)\newline \checkbox\ Modificar\newline \checkbox\ Eliminar \\
\hline
R14 & Auditoría en la CI + Dependabot con triaje semanal y alertas + casos de abuso en la Definición de Hecho. & \checkedbox\ Sí (ESLint 9.39.5 el 15-10: \texttt{pnpm audit} a cero; la US-07 cerró con 18 pruebas y sus casos de abuso; alertas de seguridad activas, ninguna abierta; Dependabot sin propuestas en la semana y sin PRs abiertas).\newline \checkbox\ No & — & \checkedbox\ Mantener\newline \checkbox\ Modificar\newline \checkbox\ Eliminar \\
\hline
\end{tabularx}

\Needspace*{30\baselineskip}
\underline{\textbf{Riesgo significativo R11 (por disparador) y moderado R12 (por cambio de medida)}}

\begin{tabularx}{\textwidth}{|C{0.9cm}|Y|Y|C{2.2cm}|L{2.4cm}|}
\hline
\multicolumn{1}{|C{0.9cm}|}{\textbf{Id}} &
\multicolumn{1}{>{\centering\arraybackslash}X|}{\textbf{Medida de mitigación}} &
\multicolumn{1}{>{\centering\arraybackslash}X|}{\textbf{¿Aplicada?}} &
\multicolumn{1}{C{2.2cm}|}{\textbf{Justificación (del NO)}} &
\multicolumn{1}{C{2.4cm}|}{\textbf{Decisión}} \\
\hline
R11 & Todo en git con remoto y subidas frecuentes. & \checkbox\ Sí\newline \checkedbox\ No, en parte: tres sesiones seguidas cerraron sin subir al remoto (11-10, 12-10 y la primera del 15-10). & La subida necesita el OK del autor y las tres sesiones terminaron sin pedirlo. Desde el 15-10, cada sesión cierra con su subida. & \checkedbox\ Mantener (la medida es correcta; falló su ejecución)\newline \checkbox\ Modificar\newline \checkbox\ Eliminar \\
\hline
R12 & Capítulos a revisión conforme se cierran; plazo orientativo por envío; entrega quincenal de la memoria. & \checkbox\ Sí\newline \checkedbox\ No, en parte: la entrega del 15-10 no se hizo. & El tutor no había respondido a la del 05-10; ese día el autor acordó con él entregas flexibles, según la disponibilidad de ambos. & \checkbox\ Mantener\newline \checkedbox\ Modificar\newline \checkbox\ Eliminar \\
\hline
\end{tabularx}

% Dos bloques: las cinco filas juntas no caben en una página (tabularx no parte).
\Needspace*{28\baselineskip}
\underline{\textbf{Cambios realizados en las medidas de mitigación (R1, R3 y R8)}} \updatesymbol

\begin{tabularx}{\textwidth}{|C{1.4cm}|Y|Y|Y|}
\hline
\multicolumn{1}{|C{1.4cm}|}{\textbf{Id}} &
\multicolumn{1}{>{\centering\arraybackslash}X|}{\textbf{Medida original}} &
\multicolumn{1}{>{\centering\arraybackslash}X|}{\textbf{Cambio realizado}} &
\multicolumn{1}{>{\centering\arraybackslash}X|}{\textbf{Justificación}} \\
\hline
R1 & Indicador: «el sprint activo gana más de 2 puntos tras la planificación del lunes…». & Indicador: «…tras la planificación semanal…». & La planificación vale cualquier día de la semana desde el 15-10 (cap.~10). \\
\hline
R3 & «…y su informe se tría en la planificación del lunes.» & «…y su informe lo aplica el asistente en la semana siguiente, preguntando al autor solo las decisiones.» & Decisión del autor del 15-10: el informe es demasiado largo para leerlo en sesiones cortas; prefiere leer la memoria y corregir a mano o pedir cambios concretos. El procedimiento (\texttt{/coherencia aplicar}) queda en los README, no en el plan. \\
\hline
R8 & Horas reales por sesión en el diario; cada retrospectiva las compara con las declaradas; calendario con holgura. & Se añade: las semanas sin disponibilidad conocidas de antemano se declaran a cero antes de empezar (cap.~10) y quedan fuera de la comparación. & Decisión del autor del 17-10: una semana a cero prevista no es capacidad perdida, y así el indicador no la cuenta como tal. \\
\hline
\end{tabularx}

\Needspace*{34\baselineskip}
\underline{\textbf{Cambios realizados en las medidas de mitigación (R10 y R12: indicador, contingencia y medida)}} \updatesymbol

\begin{tabularx}{\textwidth}{|C{1.4cm}|Y|Y|Y|}
\hline
\multicolumn{1}{|C{1.4cm}|}{\textbf{Id}} &
\multicolumn{1}{>{\centering\arraybackslash}X|}{\textbf{Medida original}} &
\multicolumn{1}{>{\centering\arraybackslash}X|}{\textbf{Cambio realizado}} &
\multicolumn{1}{>{\centering\arraybackslash}X|}{\textbf{Justificación}} \\
\hline
R10 & Indicador: «más de tres días consecutivos sin poder trabajar por causa sobrevenida dentro de un periodo de dedicación plena». Contingencia: pausa formal, re-planificación y, en el extremo, cambio de convocatoria. & Indicador: «una semana de sprint sin poder trabajar por causa sobrevenida, es decir, no declarada a cero antes de empezar (cap.~10); las semanas que reúnen sus horas solo en el fin de semana no cuentan». Contingencia con un primer escalón, el modo mínimo: si un imprevisto amenaza con dejar la semana a cero, una sesión de al menos 1~h de diario y memoria para no perder el hilo; a partir de la segunda semana consecutiva que, sin haberse declarado a cero, queda a cero o en modo mínimo, la contingencia anterior. & Lo que mide la capacidad son las horas de la semana, no los días (regla del autor del 17-10): el indicador por días saltaba con semanas que cumplían su capacidad en el fin de semana. Y la contingencia carecía de un primer escalón que evitase perder el hilo, que fue el coste real del parón de agosto a octubre. El umbral de la segunda semana consecutiva se fija en este seguimiento: una semana perdida la absorbe la holgura del plan, como una declarada a cero (cap.~10), y lo que R10 vigila es la indisponibilidad prolongada; escalar ya en la primera (alternativa descartada) convertiría cada imprevisto en una re-planificación. Las declaradas a cero no cuentan para el umbral: el cap.~10 las excluye de R10. \\
\hline
R12 & «…y cada quince días (los días 1 y 15 de cada mes) se envía al tutor la versión vigente de la memoria…»; «acordar en cada envío un plazo orientativo de respuesta». & Enviar los capítulos a revisión conforme se van cerrando y la versión vigente de la memoria cuando haya cambios que merezcan revisión, según la disponibilidad de ambos (acordado con el tutor el 15-10). Se retiran la cadencia quincenal y el plazo orientativo. & Una cadencia fija que la otra parte no puede seguir no mitiga nada. El indicador se mantiene (decisión del autor en este seguimiento): el riesgo es precisamente la demora, y suavizarlo justo antes de que salte lo vaciaría. \\
\hline
\end{tabularx}

\Needspace*{8\baselineskip}
\underline{\textbf{Nuevas medidas de mitigación a añadir}} \updatesymbol

\begin{tabularx}{\textwidth}{|C{1.4cm}|Y|Y|}
\hline
\multicolumn{1}{|C{1.4cm}|}{\textbf{Id}} &
\multicolumn{1}{>{\centering\arraybackslash}X|}{\textbf{Nueva medida}} &
\multicolumn{1}{>{\centering\arraybackslash}X|}{\textbf{Justificación}} \\
\hline
— & Ninguna en este seguimiento. & — \\
\hline
\end{tabularx}

\bigline

\section{Riesgos a revisar en este seguimiento}

Los riesgos \textbf{sin novedad y sin observación relevante} se agrupan en la tabla siguiente (indicador en su valor esperado, sin salto, estado «sin novedad»); los que dejaron señal en el periodo (R11, R12 y R14) conservan su ficha completa a continuación.

\Needspace*{18\baselineskip}
\begin{tabularx}{\textwidth}{|C{0.9cm}|Y|C{2.6cm}|}
\hline
\multicolumn{1}{|C{0.9cm}|}{\textbf{Id}} &
\multicolumn{1}{>{\centering\arraybackslash}X|}{\textbf{Indicador: valor real (esperado 0 en todos)}} &
\multicolumn{1}{C{2.6cm}|}{\textbf{Nuevo estado}} \\
\hline
R1 & 0 — el sprint no ganó puntos tras la planificación del 12-10 y nada entró fuera del backlog. El indicador se reescribe («planificación semanal»), sin cambiar su umbral. & Sin novedad \\
\hline
R2 & Indicador (tres primeros sprints) ya evaluado en S3. S8: 2 SP en un sprint cuyas horas no se reparten por tarea; la equivalencia no es medible este sprint y la decisión sigue fijada para S9. & Sin novedad \\
\hline
R3 & 0 — caps.~9 y 10 actualizados en la sesión del cambio; diario escrito en cada sesión; informe de coherencia aplicado en la semana. & Sin novedad \\
\hline
R4 & 0 — cinco fuentes nuevas, citadas al escribir. & Sin novedad \\
\hline
R5 & 0 — las decisiones del refactor constan en el diario con sus alternativas, y el autor exigió y obtuvo la justificación de cada una (la del analizador de JSON global, con la evidencia de un proyecto anterior) antes de aprobarla. & Sin novedad \\
\hline
R6 & 0 — la US-07 (2 SP, 4~h previstas) se cerró en su sprint sin desvío. & Sin novedad \\
\hline
R7 & 0 — la IP del cliente, único dato personal nuevo, no se persiste y se censura en los logs; el modelo de datos no cambió. & Sin novedad \\
\hline
R8 & 0 — 8~h~30 frente a 5 declaradas (170~\%), muy por encima del 70~\%. & Sin novedad \\
\hline
R10 & 0 — sin parón: cuatro sesiones (lunes, jueves, sábado y domingo). Su indicador y su contingencia se redefinen (arriba) sin que haya saltado. & Sin novedad \\
\hline
\end{tabularx}

% ------------------ R11 ------------------
\Needspace*{30\baselineskip}
\subsection*{Riesgo: \markfield{R11 — Pérdida de trabajo o del entorno local} (por disparador)}

\textbf{Indicadores (clave y de resolución)}

\begin{tabularx}{\textwidth}{|Z|Z|Z|Z|Z|}
\hline
\textbf{Indicador} & \textbf{Valor esperado} & \textbf{Valor real} & \textbf{¿Cumplido?} & \textbf{Acción / observaciones} \\
\hline
Reinicios espontáneos, errores de disco o arranques anómalos del equipo, o más de dos sesiones consecutivas sin subir los cambios al remoto. & 0 & \textbf{Sí, el 15-10}: tres sesiones consecutivas sin subir, la del 11-10 (commit \texttt{cebf2e8}), la del 12-10 y la primera del 15-10. La segunda sesión del 15-10 subió todo y el 17-10 hubo tres subidas con cuatro commits (la US-07 y el refactor de la API subieron juntos). Sin anomalías del equipo. & \checkedbox\ Sí\quad \checkbox\ No & Observado por el autor el 15-10, que decidió registrarlo aquí en lugar de abrir una revisión inmediata a tres días del cierre. Pérdida potencial acotada a esas tres sesiones (unas 3~h). Despejado el mismo 15-10. \\
\hline
\end{tabularx}

\textbf{Nueva probabilidad (1–3):} \underline{\ 1\ } · \textbf{Nuevo impacto (1–3):} \underline{\ 3\ } · \textbf{Nueva severidad:} \checkbox\ Crítico · \checkedbox\ Significativo · \checkbox\ Moderado\ \updatesymbol\ (sin cambio: el episodio no revela una probabilidad mayor de fallo del equipo; el fallo fue de disciplina y está corregido)

\textbf{¿Cerrar el riesgo?} \checkbox\ Sí · \checkedbox\ No\ \updatesymbol

\textbf{¿Se ha materializado?} \checkbox\ Sí · \checkedbox\ No\ \updatesymbol

\textbf{Nuevo estado} (cap.~7): \checkedbox\ Identificado (indicador saltado el 15-10 y despejado el mismo día; el plan v7 lo anota) · \checkbox\ Sin novedad · \checkbox\ Activo · \checkbox\ Materializado · \checkbox\ Resuelto · \checkbox\ Cerrado

\textbf{Plan de contingencia (solo si se ha materializado)}: no aplicable.
Causa y lección: la subida requiere el OK del autor en cada sesión (regla de CLAUDE.md) y las tres terminaron sin pedirlo; la regla es correcta, pero la sesión no debe cerrarse sin pedirlo, y desde el 15-10 forma parte del cierre de cada sesión.

% ------------------ R12 ------------------
\Needspace*{30\baselineskip}
\subsection*{Riesgo: \markfield{R12 — Ciclos de revisión con el tutor} (por cambio de medida)}

\textbf{Indicadores (clave y de resolución)}

\begin{tabularx}{\textwidth}{|Z|Z|Z|Z|Z|}
\hline
\textbf{Indicador} & \textbf{Valor esperado} & \textbf{Valor real} & \textbf{¿Cumplido?} & \textbf{Acción / observaciones} \\
\hline
Consulta/capítulo >2 semanas sin respuesta, o revisión con cambios estructurales sobre capítulo cerrado. & 0 & El correo del 05-10 (memoria completa y cuatro consultas, entre ellas la del cap.~3, §3.2 frente a §3.5.1) sigue sin respuesta a 18-10: 13 días, dentro del plazo. La entrega del 15-10 no se hizo; ese día el autor escribió al tutor y acordaron entregas flexibles. & \checkbox\ Sí\quad \checkedbox\ No & El indicador se mantiene (decisión del autor en este seguimiento): si el 19-10 sigue sin respuesta, salta y se registra en S9 con R12 activo. La medida sí cambia (ver «Cambios realizados»). Cierra la nota del S7: la consulta del cap.~3 iba en el correo del 05-10. \\
\hline
\end{tabularx}

\textbf{Nueva probabilidad (1–3):} \fieldline{0.8cm} · \textbf{Nuevo impacto (1–3):} \fieldline{0.8cm} · \textbf{Nueva severidad:} \checkbox\ Crítico · \checkbox\ Significativo · \checkbox\ Moderado\ \updatesymbol\ (sin reevaluación: el indicador no ha saltado)

\textbf{¿Cerrar el riesgo?} \checkbox\ Sí · \checkedbox\ No\ \updatesymbol

\textbf{¿Se ha materializado?} \checkbox\ Sí · \checkedbox\ No\ \updatesymbol

\textbf{Nuevo estado} (cap.~7): \checkbox\ Identificado · \checkedbox\ Sin novedad · \checkbox\ Activo · \checkbox\ Materializado · \checkbox\ Resuelto · \checkbox\ Cerrado

\textbf{Plan de contingencia (solo si se ha materializado)}: no aplicable.

% ------------------ R14 ------------------
\Needspace*{30\baselineskip}
\subsection*{Riesgo: \markfield{R14 — Verificación de seguridad insuficiente} (activo desde S7)}

\textbf{Indicadores (clave y de resolución)}

\begin{tabularx}{\textwidth}{|Z|Z|Z|Z|Z|}
\hline
\textbf{Indicador} & \textbf{Valor esperado} & \textbf{Valor real} & \textbf{¿Cumplido?} & \textbf{Acción / observaciones} \\
\hline
Historia con autenticación, autorización o datos personales cerrada sin casos de abuso, o vulnerabilidad alta sin resolver: las de producción, dentro del sprint; las de desarrollo, al cierre del siguiente. & 0 & 0 — la US-07 (autenticación) cerró con 18 pruebas y sus casos de abuso (ráfagas, bloqueo por cuenta e IP, correo malformado, cabecera de IP falsa). Las siete altas de desarrollo que activaron el riesgo en S7 se resolvieron el 15-10 con ESLint 9.39.5: \texttt{pnpm audit} a cero en producción y en desarrollo, dentro del plazo (cierre de S8). Alertas de seguridad activas, ninguna abierta; Dependabot sin propuestas en la semana. & \checkbox\ Sí\quad \checkedbox\ No & El indicador que lo activó está despejado: vuelve a identificado en el plan v7, con nota. La subida de ESLint fue manual, porque Dependabot no la había propuesto: lo que la forzó fue el plazo del indicador precisado en S7, que seguirá valiendo cuando Dependabot no proponga la corrección; la medida no cambia. \\
\hline
\end{tabularx}

\textbf{Nueva probabilidad (1–3):} \fieldline{0.8cm} · \textbf{Nuevo impacto (1–3):} \fieldline{0.8cm} · \textbf{Nueva severidad:} \checkbox\ Crítico · \checkbox\ Significativo · \checkbox\ Moderado\ \updatesymbol\ (sin reevaluación: el indicador no ha saltado este sprint; el salto de S7 se despeja)

\textbf{¿Cerrar el riesgo?} \checkbox\ Sí · \checkedbox\ No\ \updatesymbol

\textbf{¿Se ha materializado?} \checkbox\ Sí · \checkedbox\ No\ \updatesymbol

\textbf{Nuevo estado} (cap.~7): \checkedbox\ Identificado (activo en S7; indicador despejado el 15-10) · \checkbox\ Sin novedad · \checkbox\ Activo · \checkbox\ Materializado · \checkbox\ Resuelto · \checkbox\ Cerrado

\textbf{Plan de contingencia (solo si se ha materializado)}: no aplicable.

\Needspace*{10\baselineskip}
\bigline

\section{Reflexión desde el último seguimiento}

\textbf{¿Ha ocurrido algo reseñable desde el último seguimiento?} \checkedbox\ Sí · \checkbox\ No

Una semana desde el seguimiento del S7.
La US-07 cerró EP02, y con ella llegó el refactor de la API por funcionalidad (17-10).
El método cambió en tres puntos: la planificación vale cualquier día de la semana (15-10), el informe de coherencia lo aplica el asistente y el autor solo decide (15-10), y la semana a cero entra en el cap.~10 y el modo mínimo en el README del diario (17-10); los tres se trasladan al plan en esta revisión.
El seguimiento del S7 y la v6 se corrigieron en los propios documentos (errores objetivos el 15-10; el 17-10, R14 pasó a activo y se añadió la línea del plan emitido, que esta plantilla incorpora).
La entrega quincenal al tutor se retiró el 15-10.

\Needspace*{6\baselineskip}
\textbf{Nuevos riesgos identificados} \updatesymbol

\begin{tabularx}{\textwidth}{|C{1.4cm}|Y|C{0.8cm}|C{0.8cm}|C{1cm}|C{2.4cm}|}
\hline
\multicolumn{1}{|C{1.4cm}|}{\textbf{Id}} &
\multicolumn{1}{>{\centering\arraybackslash}X|}{\textbf{Descripción}} &
\multicolumn{1}{C{0.8cm}|}{\textbf{P}} &
\multicolumn{1}{C{0.8cm}|}{\textbf{I}} &
\multicolumn{1}{C{1cm}|}{\textbf{Exp.}} &
\multicolumn{1}{C{2.4cm}|}{\textbf{Severidad}} \\
\hline
— & Ninguno. & — & — & — & — \\
\hline
\end{tabularx}

\textbf{¿Afecta algo a los objetivos o al alcance (cap.~3 de la memoria)?} \checkbox\ Sí · \checkedbox\ No

\textbf{Justificación:} EP02 queda completa (US-04 a US-07) y el alcance no cambia; a F2 le quedan 34 SP para 14 sprints (35 SP de capacidad a 2,0~h/SP): sobra un punto de capacidad, cuando en S7 sobraba un punto de trabajo (41 SP frente a 40), porque los dos primeros sprints entregaron 7 SP con 10~h de capacidad (5 SP).

\textbf{Otros eventos relevantes:} CLAUDE.md recortado de 168 a 101 líneas (15-10), con lo operativo repartido en los README de cada carpeta; informe de coherencia del 11-10 aplicado el 15-10 (23 hallazgos: 22 aplicados y 1 que ya no aplicaba); refactor de la API por funcionalidad, con sus pruebas en espejo (17-10); GitHub anuncia que \texttt{ubuntu-latest} pasa a Ubuntu~26 desde el 19-10 (aviso en la CI del 17-10): vigilar la primera ejecución posterior (R13); rama local \texttt{temp-backup} borrada; Dependabot sin propuestas y sin PRs abiertas; límite de sesión de la suscripción agotado el 15-10 a mitad de una revisión con tres flujos de agentes en paralelo (unos 9 millones de tokens): no es el indicador de R15, que mide la cuota semanal y las tareas recurrentes dimensionadas, pero es su primer dato y se revisa con él en S9.

\end{document}
